\documentclass{zhvt-classic}

% 使用本地源流明体（GenRyuMin2 TW）按文件路径加载，不依赖系统字体名
% 版框（内框／界栏／书口侧栏／鱼尾）改用与版心印章同一印泥棕色：
%   印章图片 yshs.png 实色 = RGB(135,68,52) = HTML 874434
\usepackage{xcolor}
\usepackage{tikz}
% sealbrown 双注册：xcolor 供 TikZ/pgf，L3 供 cls 版框
\definecolor{sealbrown}{HTML}{874434}
\ExplSyntaxOn
\color_set:nnn { sealbrown } { HTML } { 874434 }
\ExplSyntaxOff

\zhvtset{
  grid color       = sealbrown,
  main font        = GenRyuMin2TW-SB,
  main font path   = C:/Users/Administrator/texmf/fonts/opentype/public/GenRyuMin2TW/,
jiazhu font      = GenRyuMin2TW-M,
  jiazhu font path = C:/Users/Administrator/texmf/fonts/opentype/public/GenRyuMin2TW/,
  % 页码框容纳 4 字：正常页「九十七」+前置空 1 格=4 行恰好；
  % 100+ 页「一百零一」共 5 行超出，末字「一/二/三」被框裁掉（2026-10-04 实测），故放宽到 5
  shukou page number chars = 5,
  % 书口书名框高从默认 6 字减到 4 字，把省出的 2 字高拨给 chapter 框，
  % 使 6 字书名+卷名放进 9 字高的 chapter 框后不再压线（2026-10-05）
  shukou title chars = 4,
}

\title{史记}[黥布列传]
\author{汉·司马迁}
\maker{闲叟漫记·重製}


% ===== 索隐述赞样式（2026-11-08 四次定稿 · 愚大人亲手改定参数）=====
% ⚠ 竖排前提：自定义字体家族必须带 RawFeature={vertical:+vert}；
%   simhei.ttf 加 +vhal 会整体加载失败静默回退，只写 +vert；simfang 同理。
\newCJKfontfamily{\shuzanfont}[Path=C:/Windows/Fonts/,RawFeature={vertical:+vert}]{simhei.ttf}
\newCJKfontfamily{\shuzanbody}[Path=C:/Windows/Fonts/,RawFeature={vertical:+vert}]{simfang.ttf}
% 题头：152x30pt 圆角 6pt 棕底白字；baseline (0,2) 使整框落入界行格子不压版框线；
% 白字节点 x 不补偿、y +8.2pt（像素级墨迹实测 dx/dy≈0，L/R 留白约 20pt）。
\newcommand{\shuzanhead}[1]{%
  \begin{tikzpicture}[baseline={(0,2)}]
    \fill[sealbrown, rounded corners=6pt]
      (-76pt,-15pt) rectangle (76pt,15pt);
    \node[text=white, inner sep=0pt, anchor=center] at (0pt,8.2pt)
      {\shuzanfont\fontsize{23}{23}\selectfont #1};
  \end{tikzpicture}%
}
% ★ 修正书口卷名位置：chapter 框起始位置改回 cls 默认值，避免全局影响。
%   只在每个卷首页通过 \juan 命令里的空命令微调该页位置（安全，不会跨页污染）。
\ExplSyntaxOn
\AtBeginDocument{
  \dim_set:Nn \l__zhvt_shukou_chapter_title_start_dim
    { \l__zhvt_shukou_title_height_dim + \l__zhvt_font_size_dim * 2 }
  \dim_set:Nn \l__zhvt_shukou_page_number_start_dim
    { \l__zhvt_shukou_chapter_title_start_dim + \l__zhvt_shukou_chapter_title_height_dim }
}
\ExplSyntaxOff

\begin{document}

% ★★ 版心（书口）排布覆写：把书名放到鱼尾之下（2026-10-04）
%
%   cls 默认版心自上而下：书名（title 静态框，在鱼尾之上）→ 鱼尾 → 篇名（chapter 动态框）→ 页码
%   本模板改为：          鱼尾（位置不动）→ 书名 → 篇名 → 页码分割线 → 页码
%
%   手法：把 cls 放在 title 静态框里的书名，改放进 chapter 动态框（鱼尾之下）；
%         title 静态框留空。书名取 \l__zhvt_title_str，即 \title{...} 的第一参数，
%         所以换书时只改 \title，不必在此重复写书名。
%
%   ⚠ 三条必须遵守的写法（2026-10-04 实测）：
%     ① 汉字必须用 \hbox{...} 包住。flowfram 会把框内容装进一个宽 14.4mm（= 一个 line_sep
%        = 书口栏宽）的纵向 minipage 并按该宽度折行，而整页靠 PDF /Rotate 90 显示、
%        框内文字本身不变向 —— 不包 \hbox 时「围炉夜话」4 字会被逐字折成 4 行，
%        在版面上就成了横排 4 字并横向溢出版框。
%     ② 竖排下**一个 \hbox 就是一整列**。书名与篇名要同处一列，必须放进同一个 \hbox；
%        分成两个 \hbox 时，后一个会被当成新的一列、甩到隔壁栏去
%        （实测「序」跑到书口栏外侧 14.4mm 的正文列上）。中间的 \quad 在框内是横向空白，
%        显示出来正好是书名与篇名之间的纵向间隔。
%     ③ ★ 内容盒要写成 **与框等宽的定宽盒** `\hbox to \l__zhvt_shukou_chapter_title_height_dim{…\hfil}`。
%        实测：chapter 框是把内容盒沿纵向**居中**摆放的，内容宽度一变，书名就跟着漂：
%        自然宽度的 \hbox（书名 38.4mm）落在 x=157.2mm，加篇名后（57.6mm）落在 x=152.4mm，
%        目录页（无篇名）与正文页的书名就差半格。做成等宽盒后，无论有没有篇名，
%        内容固定从框顶起排，实测两页均落在 x=142.8mm，不再漂。
\ExplSyntaxOn
% 把书口 chapter 框（书名+篇名）起始位置下移 1 字高，避免 4 字卷名顶到鱼尾
% ⚠ cls 的 \__zhvt_calculate_dimensions: 在 \AtBeginDocument 里执行，会重算该 dim；
%   必须把覆写挂到 begin document 之后，否则会被二次覆盖（2026-10-05 实测）。
% (旧方案已移到导言区，直接修改 chapter 框起始位置)
% 卷标题命令 #1=书口篇名(鱼尾) #2=正文卷名（含卷序，卷名内空格用 ~）
%
% ★★ 目录页两项修订（2026-10-05，愚大人指出）★★
%   ① 页码起点统一对齐：
%      原状 —— 目录项写成 \par\hspace{4em}卷名\hspace{4em}页码，页码起点 = 卷名实宽 + 4em。
%      卷名实宽随字数变：「卷一~梁惠王上」6 字 +~、「卷七~离娄上」5 字 +~，
%      差 1 个字位 → 卷七~卷十 4 条的页码整体比其余 10 条高 27.21pt（一个字），看着上下不齐。
%      改法 —— 卷名补全角空格（宽 = 1 字 = 27.215pt）补足到 **7 字符位**再写进目录，
%      14 条页码起点遂全部落在同一条竖线上（起点对齐；末字随页码字数自然参差）。
%      ⚠ 只对目录里的副本 \l_tmpc_tl 补，正文卷名仍用 \l_tmpa_tl 原样 —— 正文像素不受影响。
%   ② 卷名整体上移半字位：前导 \hspace 由 4em 改 3.5em。
%      实测原状目录块在版框内偏下（天头留白 139.6pt / 地脚留白 110.7pt），
%      上移 0.5em = 13.6pt 后上下基本均分，视觉居中。
\NewDocumentCommand { \juan } { mm } {
	\clearpage
	\tl_set:Nn \l__zhvt_chapter_title_tl { #1 }
	\tl_set:Nn \l_tmpa_tl { #2 }
	\tl_set_eq:NN \l_tmpc_tl \l_tmpa_tl
	\int_step_inline:nn { 7 - \str_count:n { #2 } }
	  { \tl_put_right:Nn \l_tmpc_tl { 　 } }
	\tl_set:Nx \l_tmpb_tl { \zhnum { page } }
	\tl_put_left:Nn \l_tmpb_tl { \hspace { 4em } }
	\tl_put_left:Nn \l_tmpb_tl { \l_tmpc_tl }
	\tl_put_left:Nn \l_tmpb_tl { \par \hspace { 3.5em } }
	\addtocontents { toc } { \l_tmpb_tl }
	\skip_horizontal:n { \l__zhvt_font_size_dim * 4 }
	\l_tmpa_tl
	\par
	\thispagestyle { empty }
}
\cs_set_protected:Npn \__zhvt_set_shukou_content: {
	% ★ 鱼尾上方留白（title 静态框）嵌入闲叟钤印：
	%   整页 /Rotate 90，故此处把图用 \rotatebox{90} 转正为竖章；
	%   尺寸卡在版心条幅内（书内约 7mm 宽 × 33.8mm 高，留余量不压鱼尾）。
	\setstaticcontents * { title } {
		\vbox_to_ht:nn { \l__zhvt_baseline_skip_dim } {
			\vskip 14.2pt
			\hbox_to_wd:nn { \l__zhvt_shukou_title_height_dim } {
				\hfil
				\rotatebox [ origin = c ] { 90 } { \includegraphics [ width = 7mm , height = 33.8mm ] { yshs.png } }
				\hfil
			}
		}
	}
% ★ 版心（书口）统一改法（2026-11-16 愚大人定）：18pt＋框内居中＋书口文字同版框色
%   ① 原 27.2pt 时 7 字长篇名（如「廉颇蔺相如列传」）末字溢出 chapter 框、压住页码；
%      改 \fontsize{18pt}{18pt}，前后 \hfil 令字串在框内居中（不随篇名字数漂移）。
%   ② \vbox{\vskip 8.2pt} 定量下移，抵消缩字号后栏宽方向 8.2pt 上偏（像素实测 y 中心 543.3 = 原版）。
%      ⚠ \raisebox / \vspace* 在 flowfram 框内无效；换字号须同步换补偿：16pt→10.1、19pt→7.4。
%   ③ chapter 与 page~number 均置 \color{sealbrown}（HTML 874434），与版框／界行同色；正文仍黑。
%   ⚠ 页码字位、钤印、鱼尾、界行、正文一律未动；五部合订本同此口径。
	\setdynamiccontents * { chapter }
	{ \color{sealbrown} \vbox { \vskip 8.2pt \hbox to \l__zhvt_shukou_chapter_title_height_dim
		{ \hfil \fontsize{18pt}{18pt}\selectfont \l__zhvt_title_str \quad \l__zhvt_chapter_title_tl \hfil } } }
	\setdynamiccontents * { page~number } { \color{sealbrown} \quad \zhnum { page } }
}
\__zhvt_set_shukou_content:
\ExplSyntaxOff

\maketitle

% 单篇，无多卷目录，不渲染目录页

\juan{黥布列传}{卷九十一~黥布列传第三十一}
\hspace*{2em}黥布者。六人也。姓英氏。秦时为布衣。少年。有客相之曰。﹁当刑而王。﹂及壮。坐法黥。布欣然笑曰。﹁人相我当刑而王。几是乎。﹂人有闻者。共俳笑之。布已论输丽山。丽山之徒数十万人。布皆与其徒长豪桀交通。乃率其曹偶。亡之江中为群盗。\\
\hspace*{2em}陈胜之起也。布乃见番君。与其众叛秦。聚兵数千人。番君以其女妻之。章邯之灭陈胜。破吕臣军。布乃引兵北击秦左右校。破之清波。引兵而东。闻项梁定江东会稽。涉江而西。陈婴以项氏世为楚将。乃以兵属项梁。渡淮南。英布。蒲将军亦以兵属项梁。\\
\hspace*{2em}项梁涉淮而西。击景驹。秦嘉等。布常冠军。项梁至薛。闻陈王定死。乃立楚怀王。项梁号为武信君。英布为当阳君。项梁败死定陶。怀王徙都彭城。诸将英布亦皆保聚彭城。当是时。秦急围赵。赵数使人请救。怀王使宋义为上将。范曾为末将。项籍为次将。英布。蒲将军皆为将军。悉属宋义。北救赵。及项籍杀宋义於河上。怀王因立籍为上将军。诸将皆属项籍。项籍使布先渡河击秦。布数有利。籍乃悉引兵涉河从之。遂破秦军。降章邯等。楚兵常胜。功冠诸侯。诸侯兵皆以服属楚者。以布数以少败众也。\\
\hspace*{2em}项籍之引兵西至新安。又使布等夜击阬章邯秦卒二十馀万人。至关。不得入。又使布等先从间道破关下军。遂得入。至咸阳。布常为军锋。项王封诸将。立布为九江王。都六。\\
\hspace*{2em}汉元年四月。诸侯皆罢戏下。各就国。项氏立怀王为义帝。徙都长沙。乃阴令九江王布等行击之。其八月。布使将击义帝。追杀之郴县。\\
\hspace*{2em}汉二年。齐王田荣畔楚。项王往击齐。徵兵九江。九江王布称病不往。遣将将数千人行。汉之败楚彭城。布又称病不佐楚。项王由此怨布。数使使者诮让召布。布愈恐。不敢往。项王方北忧齐。赵。西患汉。所与者独九江王。又多布材。欲亲用之。以故未击。\\
\hspace*{2em}汉三年。汉王击楚。大战彭城。不利。出梁地。至虞。谓左右曰。﹁如彼等者。无足与计天下事。﹂谒者随何进曰。﹁不审陛下所谓。﹂汉王曰。﹁孰能为我使淮南。令之发兵倍楚。留项王於齐数月。我之取天下可以百全。﹂随何曰。﹁臣请使之。﹂乃与二十人俱。使淮南。至。因太宰主之。三日不得见。随何因说太宰曰。﹁王之不见何。必以楚为彊。以汉为弱。此臣之所以为使。使何得见。言之而是邪。是大王所欲闻也。言之而非邪。使何等二十人伏斧质淮南市。以明王倍汉而与楚也。﹂太宰乃言之王。王见之。随何曰。﹁汉王使臣敬进书大王御者。窃怪大王与楚何亲也。﹂淮南王曰。﹁寡人北乡而臣事之。﹂随何曰。﹁大王与项王俱列为诸侯。北乡而臣事之。必以楚为彊。可以讬国也。项王伐齐。身负板筑。以为士卒先。大王宜悉淮南之众。身自将之。为楚军前锋。今乃发四千人以助楚。夫北面而臣事人者。固若是乎。夫汉王战於彭城。项王未出齐也。大王宜骚淮南之兵渡淮。日夜会战彭城下。大王抚万人之众。无一人渡淮者。垂拱而观其孰胜。夫讬国於人者。固若是乎。大王提空名以乡楚。而欲厚自讬。臣窃为大王不取也。然而大王不背楚者。以汉为弱也。夫楚兵虽彊。天下负之以不义之名。以其背盟约而杀义帝也。然而楚王恃战胜自彊。汉王收诸侯。还守成皋。荥阳。下蜀。汉之粟。深沟壁垒。分卒守徼乘塞。楚人还兵。间以梁地。深入敌国八九百里。欲战则不得。攻城则力不能。老弱转粮千里之外。楚兵至荥阳。成皋。汉坚守而不动。进则不得攻。退则不得解。故曰楚兵不足恃也。使楚胜汉。则诸侯自危惧而相救。夫楚之彊。適足以致天下之兵耳。故楚不如汉。其势易见也。今大王不与万全之汉而自讬於危亡之楚。臣窃为大王惑之。臣非以淮南之兵足以亡楚也。夫大王发兵而倍楚。项王必留。留数月。汉之取天下可以万全。臣请与大王提剑而归汉。汉王必裂地而封大王。又况淮南。淮南必大王有也。故汉王敬使使臣进愚计。原大王之留意也。﹂淮南王曰。﹁请奉命。﹂阴许畔楚与汉。未敢泄也。\\
\hspace*{2em}楚使者在。方急责英布发兵。舍传舍。随何直入。坐楚使者上坐。曰。﹁九江王已归汉。楚何以得发兵。﹂布愕然。楚使者起。何因说布曰。﹁事已搆。可遂杀楚使者。无使归。而疾走汉并力。﹂布曰。﹁如使者教。因起兵而击之耳。﹂於是杀使者。因起兵而攻楚。楚使项声。龙且攻淮南。项王留而攻下邑。数月。龙且击淮南。破布军。布欲引兵走汉。恐楚王杀之。故间行与何俱归汉。\\
\hspace*{2em}淮南王至。上方踞床洗。召布入见。布大怒。悔来。欲自杀。出就舍。帐御饮食从官如汉王居。布又大喜过望。於是乃使人入九江。楚已使项伯收九江兵。尽杀布妻子。布使者颇得故人幸臣。将众数千人归汉。汉益分布兵而与俱北。收兵至成皋。四年七月。立布为淮南王。与击项籍。\\
\hspace*{2em}汉五年。布使人入九江。得数县。六年。布与刘贾入九江。诱大司马周殷。周殷反楚。遂举九江兵与汉击楚。破之垓下。\\
\hspace*{2em}项籍死。天下定。上置酒。上折随何之功。谓何为腐儒。为天下安用腐儒。随何跪曰。﹁夫陛下引兵攻彭城。楚王未去齐也。陛下发步卒五万人。骑五千。能以取淮南乎。﹂上曰。﹁不能。﹂随何曰。﹁陛下使何与二十人使淮南。至。如陛下之意。是何之功贤於步卒五万人骑五千也。然而陛下谓何腐儒。为天下安用腐儒。何也。﹂上曰。﹁吾方图子之功。﹂乃以随何为护军中尉。布遂剖符为淮南王。都六。九江。庐江。衡山。豫章郡皆属布。\\
\hspace*{2em}七年。朝陈。八年。朝雒阳。九年。朝长安。\\
\hspace*{2em}十一年。高后诛淮阴侯。布因心恐。夏。汉诛梁王彭越。醢之。盛其醢遍赐诸侯。至淮南。淮南王方猎。见醢。因大恐。阴令人部聚兵。候伺旁郡警急。\\
\hspace*{2em}布所幸姬疾。请就医。医家与中大夫贲赫对门。姬数如医家。贲赫自以为侍中。乃厚餽遗。从姬饮医家。姬侍王。从容语次。誉赫长者也。王怒曰。﹁汝安从知之。﹂具说状。王疑其与乱。赫恐。称病。王愈怒。欲捕赫。赫言变事。乘传诣长安。布使人追。不及。赫至。上变。言布谋反有端。可先未发诛也。上读其书。语萧相国。相国曰。﹁布不宜有此。恐仇怨妄诬之。请击赫。使人微验淮南王。﹂淮南王布见赫以罪亡。上变。固已疑其言国阴事。汉使又来。颇有所验。遂族赫家。发兵反。反书闻。上乃赦贲赫。以为将军。\\
\hspace*{2em}上召诸将问曰。﹁布反。为之柰何。﹂皆曰。﹁发兵击之。阬竖子耳。何能为乎。﹂汝阴侯滕公召故楚令尹问之。令尹曰。﹁是故当反。﹂滕公曰。﹁上裂地而王之。疏爵而贵之。南面而立万乘之主。其反何也。﹂令尹曰。﹁往年杀彭越。前年杀韩信。此三人者。同功一体之人也。自疑祸及身。故反耳。﹂滕公言之上曰。﹁臣客故楚令尹薛公者。其人有筹筴之计。可问。﹂上乃召见问薛公。薛公对曰。﹁布反不足怪也。使布出於上计。山东非汉之有也。出於中计。胜败之数未可知也。出於下计。陛下安枕而卧矣。﹂上曰。﹁何谓上计。﹂令尹对曰。﹁东取吴。西取楚。并齐取鲁。传檄燕。赵。固守其所。山东非汉之有也。﹂﹁何谓中计。﹂﹁东取吴。西取楚。并韩取魏。据敖庾之粟。塞成皋之口。胜败之数未可知也。﹂﹁何谓下计。﹂﹁东取吴。西取下蔡。归重於越。身归长沙。陛下安枕而卧。汉无事矣。﹂上曰。﹁是计将安出。﹂令尹对曰。﹁出下计。﹂上曰。﹁何谓废上中计而出下计。﹂令尹曰。﹁布故丽山之徒也。自致万乘之主。此皆为身。不顾後为百姓万世虑者也。故曰出下计。﹂上曰。﹁善。﹂封薛公千户。乃立皇子长为淮南王。上遂发兵自将东击布。\\
\hspace*{2em}布之初反。谓其将曰。﹁上老矣。厌兵。必不能来。使诸将。诸将独患淮阴。彭越。今皆已死。馀不足畏也。﹂故遂反。果如薛公筹之。东击荆。荆王刘贾走死富陵。尽劫其兵。渡淮击楚。楚发兵与战徐。僮间。为三军。欲以相救为奇。或说楚将曰。﹁布善用兵。民素畏之。且兵法。诸侯战其地为散地。今别为三。彼败吾一军。馀皆走。安能相救。﹂不听。布果破其一军。其二军散走。\\
\hspace*{2em}遂西。与上兵遇蕲西。会甀。布兵精甚。上乃壁庸城。望布军置陈如项籍军。上恶之。与布相望见。遥谓布曰。﹁何苦而反。﹂布曰。﹁欲为帝耳。﹂上怒骂之。遂大战。布军败走。渡淮。数止战。不利。与百馀人走江南。布故与番君婚。以故长沙哀王使人绐布。伪与亡。诱走越。故信而随之番阳。番阳人杀布兹乡民田舍。遂灭黥布。\\
\hspace*{2em}立皇子长为淮南王。封贲赫为期思侯。诸将率多以功封者。\\
\hspace*{2em}太史公曰。英布者。其先岂春秋所见楚灭英。六。皋陶之後哉。身被刑法。何其拔兴之暴也。项氏之所阬杀人以千万数。而布常为首虐。功冠诸侯。用此得王。亦不免於身为世大僇。祸之兴自爱姬殖。妒媢生患。竟以灭国。\\

% ===== 索隐述赞（司马贞《史记索隐》）· 统一样式 =====
% 题头「索隐述赞曰」五字并入反白框；正文每四句一列，定宽盒 \hfil 居列中；
% \fontsize{24}{41} 列距对齐全书界行（121+41k）；\clearpage 使述赞独占一页。
\clearpage
\shuzanhead{索隐述赞曰}\linebreak
{\shuzanbody\fontsize{24}{41}\selectfont
\hbox to \hsize{\hfil 九江初筮。当刑而王。既免徒中。聚盗江上。\hfil}\linebreak
\hbox to \hsize{\hfil 再雄楚卒。频破秦将。病为羽疑。归受汉杖。\hfil}\linebreak
\hbox to \hsize{\hfil 贲赫见毁。卒致无妄。\hfil}%
\par}
\end{document}
